% !TeX root = ../../tesis.tex

% =============================================================================
%  4.2.7 — SALIDAS
% =============================================================================
\section{Salidas: API REST con FastAPI}
\label{sec:vft-salidas}

Los resultados de VFTModel se entregan a través de una
\textbf{\api{} \textsc{rest}} construida con \textbf{FastAPI}, expuesta
localmente en \texttt{localhost:8000} con documentación interactiva disponible
en el endpoint \texttt{/docs}.

La \api{} expone cinco endpoints bajo \texttt{/api/v1/network/}:
\texttt{build-auto} (POST, construye el grafo en caché),
\texttt{spatial-coverage} (GET, cobertura por alcaldía),
\texttt{topological/capillary-strength} (GET, ranking de fuerza capilar),
\texttt{topological/geo-capillary} (GET, macro-hubs) y
\texttt{topological/detour-factor} (GET, factor de desviación por par O-D). La
especificación completa de los endpoints se presenta en el Cuadro~
\ref{tab:vft-endpoints} del Anexo~\ref{anx:vft-api}.

Cada respuesta \textsc{json} puede contener tres tipos de datos:

\begin{itemize}
    \item \textbf{Tablas de métricas:} listas de nodos con su puntuación en el
    indicador calculado, ordenadas de mayor a menor relevancia.
    \item \textbf{Geometrías de ruta:} la trayectoria real sobre la red y la
    línea recta teórica entre origen y destino, ambas en formato \geojson{},
    listas para ser renderizadas en visores cartográficos.
    \item \textbf{Resúmenes estadísticos:} promedio, máximo y distribución del
    indicador sobre el conjunto de pares O-D analizados en el muestreo.
\end{itemize}
